\section{Optimizer}
\label{sec:optimizer}

The optimizer transforms logical query plans into efficient physical
execution plans using cost-based rules.

\subsection{Cost Model}
\label{sec:optimizer:cost}

Each operator has an estimated cost based on:
\begin{itemize}
    \item Input cardinality (from table statistics or sampling).
    \item Output cardinality (from selectivity estimation).
    \item Index availability (which indexes can accelerate the operator).
    \item Engine-specific constants (sequential scan vs.\ index lookup cost).
\end{itemize}

\subsection{Optimization Rules}
\label{sec:optimizer:rules}

The optimizer applies the following rules in a fixed-point loop:

\begin{description}
    \item[Filter Pushdown] Move $\sigma$ (filter) below $\bowtie$ (join)
          to reduce input cardinality early.
    \item[Projection Pushdown] Move $\pi$ (projection) below operators to
          discard unused columns.
    \item[Join Reordering] Reorder joins by increasing estimated output
          cardinality using greedy enumeration.
    \item[Index Selection] Replace full scans with index seeks when
          applicable filters exist.
    \item[Constant Folding] Evaluate constant expressions at plan time.
\end{description}

\subsection{Engine-Specific Optimization}
\label{sec:optimizer:engine}

The KG optimizer additionally considers:
\begin{itemize}
    \item Which of the six SPO/POS/OSP permutations to use.
    \item Whether reified facts should be reconstructed via join or
          materialized view.
\end{itemize}

The Sheaf optimizer additionally considers:
\begin{itemize}
    \item Context pruning: eliminate irrelevant context subtrees.
    \item Restriction map caching: cache frequently used restrictions.
    \item Global section materialization: pre-glue known compatible sections.
\end{itemize}
